\begin{frame}[plain]
    \begin{center}
        \LARGE Fundamentação teórica
    \end{center}
\end{frame}

% \begin{frame}{Redes neurais convolucionais}
    
% \end{frame}

\begin{frame}{\textit{Convolution layer} e \textit{feature map}}
    Definição de convolução \cite{computers12080151}
\begin{equation}
    (f*g)[n]=\sum^{\infty}_{m=-\infty}f[m]g[n-m]
\end{equation}

As duas funções $f$ e $g$ podem ser discretas ou continuas, $n$ é a posição ou tempo. Considerando que os sinais de entrada sejam discretas:
\begin{equation}
    (f*g)[n]=\sum^{\infty}_{m=-\infty}f[m]g[n-m]\Delta M
\end{equation}
Onde $\Delta M$ é o intervalo de amostragem
\end{frame}


% \begin{frame}{\textit{Convolution layer} e \textit{feature map}}
%     Após o processo de convolução pela imagem, o resultado é denominado de \textit{feature maps}, representado na figura \ref{fig:featuremap} \cite{computers12080151}.
% \end{frame}

\begin{frame}[plain]
    \begin{figure}
        \centering
        \caption{Representação das \textit{convolution layer} e \textit{feature map}}
        \includegraphics[width=\linewidth]{images/theoretical-foundation/featuremap.jpg}
        \begin{center}
        {\footnotesize Fonte: Baeldung, 2023.}
    \end{center} 
        \label{fig:featuremap}
    \end{figure}
\end{frame}

% \begin{frame}{Pooling layer}
    
% \end{frame}

\begin{frame}[plain]
    \begin{figure}
        \centering
        \caption{\textit{Average pooling layer}}
        \includegraphics[width=\linewidth]{images/theoretical-foundation/Average Pooling.jpg}
        \begin{center}
        {\footnotesize Fonte: \textcite{AbhishekJain}}
    \end{center} 
        \label{fig:enter-label}
    \end{figure}
\end{frame}

\begin{frame}[plain]
    \begin{figure}
        \centering
        \caption{\textit{Max pooling}}
        \includegraphics[width=\linewidth]{images/theoretical-foundation/Max Pooling.jpg}
        \begin{center}
        {\footnotesize Fonte: \textcite{AbhishekJain}}
    \end{center} 
        \label{fig:enter-label}
    \end{figure}
\end{frame}

% \begin{frame}{\textit{Activation layer}}
%     Essa camada é responsável por adicionar uma função de não linearidade, isso faz reduzir a sua saturação e o aprendizado de características mais complexas, geralmente é adicionado ao final das camadas, seja de \textit{pooling} ou \textit{convolution} \cite{computers12080151,201811.0546}. 
% \end{frame}

\begin{frame}{Rectified linear unit (\textit{ReLU})}
    A função de ativação \textit{Rectified linear unit} (ou \textit{ReLU}) ocorre somente quando os valores são superiores a $0$ \cite{201811.0546}.

\begin{equation}\label{eq:ReLU1}
    g(y)\begin{cases}
        0\text{ para } y<0\\
        y\text{ para } y\geq 0
    \end{cases}
\end{equation}

Também pode-se representar a equação \ref{eq:ReLU1} dessa forma \cite{computers12080151}:

\begin{equation}\label{eq:ReLU2}
    f(x)=max(0,x)
\end{equation}

% As funções de ativação demonstradas nas equações \ref{eq:ReLU1} ou \ref{eq:ReLU2} não permitem valores negativos, o que o torna bom para extração de características complexas \cite{computers12080151,8286426}.
\end{frame}


\begin{frame}{\textit{Softmax}}
    A função de ativação \textit{Softmax} é empregado ao final da rede neural, para termos a distribuição de probabilidade da classe que pertence a imagem de entrada \cite{201811.0546}
\begin{equation}
    \alpha(c)_{j}=\cfrac{e^{c_{j}}}{\sum^{K}_{k=1}e^{c_{k}}}    
\end{equation}
\end{frame}

% \begin{frame}{\textit{Flatten}}
%     Processo onde utiliza a matriz de saída das operações de convolução, \textit{pooling} e entre outros para introduzir na rede neural. Para isso é necessário transformar essa matriz bidimensional em um vetor unidimensional.
% \end{frame}

\begin{frame}[plain]
    \begin{figure}
        \centering
        \caption{Processo de \textit{flatten} (ou achatamento da matriz)}
        \includegraphics[width=\linewidth]{images/theoretical-foundation/flattening.png}
        \begin{center}
        {\footnotesize Fonte: \href{https://www.superdatascience.com/blogs/convolutional-neural-networks-cnn-step-3-flattening}{SuperDataScience}, 2018.}
    \end{center} 
        \label{fig:enter-label}
    \end{figure}
\end{frame}
% \begin{frame}{\textit{Fully connected layer}}
% \end{frame}
\begin{frame}[plain]
    \begin{figure}
        \centering
        \caption{Representação dos pesos, conexões e \textit{bias} entre neurônios}
        \includegraphics[width=0.5\linewidth]{images/theoretical-foundation/neuron.png}
        \begin{center}
        {\footnotesize Fonte: Melcher, 2021}
    \end{center} 
        \label{fig:enter-label}
    \end{figure}
\end{frame}

% \begin{frame}{\textit{Fully connected layer}}
% \end{frame}

\begin{frame}[plain]
    \begin{figure}
        \centering
        \caption{Representação de uma rede neural com \textit{hidden layers}}
        \includegraphics[width=0.5\linewidth]{images/theoretical-foundation/nn.png}
        \begin{center}
        {\footnotesize Fonte: Melcher, 2021}
    \end{center} 
        \label{fig:enter-label}
    \end{figure}
\end{frame}

\begin{frame}{Processo de aprendizado}
    As arquiteturas de \textit{CNN} precisam passar pelo processo de treinamento para a definição e ajuste dos pesos em cada conexão entre os neurônios e do \textit{bias} da camada \textit{fully connected}. 
\end{frame}

\begin{frame}{\textit{Backpropagation}}
\textit{Backpropagation} é o processo de atualização e correção dos pesos de cada entrada dos neurônios e do \textit{bias} com o gradiente da função de perda, isso também pode ocorrer nas camadas de convolução se possuírem \textit {bias} \cite{computers12080151}.

\end{frame}

\begin{frame}[plain]
    \begin{align}
        \dfrac{\partial L}{\partial w^{(l)}_{i,j}}=&\dfrac{\partial L}{\partial z_{i}^{(l)}}\cdot\dfrac{\partial z_{i}^{(l)}}{\partial w_{i,j}^{(l)}}\\
        \dfrac{\partial L}{\partial b_{i}^{(l)}}=&\dfrac{\partial L}{\partial z_{i}^{(l)}}\cdot\dfrac{\partial z_{i}^{(l)}}{\partial b_{i}^{(l)}}
\end{align}

Onde:

\begin{itemize}
    \item $L$ é a função de \textit{loss} (ou perda)
    \item $w_{i,j}^{(l)}$ é o peso da conexão do neurônio $i$ na camada $l-1$ para o neurônio $j$ na camada $l$.
    \item $b_{i}^{(l)}$ é o \textit{bias} ou viés do neurônio $i$ na camada $l$.
    \item $z_{i}^{(l)}$ é a somatória das entradas do neurônio $i$ na camada $l$.
\end{itemize}
\end{frame}

\begin{frame}{\textit{Função de perda}}
    A função para calcular a perda que será utilizado será o \textit{sparse categorical cross-entropy}, implementação do \textit{Tensorflow Keras} \cite{chollet2015keras}.

    \vspace{1em}
    Utiliza a \textit{cross-entropy} (entropia cruzada) que é calculado utilizando a diferença da saída da previsão que a \textit{CNN} realizou e a classe real que a imagem pertence \cite{paszke2019pytorchimperativestylehighperformance}.
\end{frame}

\begin{frame}{Otimizador}
    O otimizador serve para encontrarmos o mínimo ou o máximo de uma função com diversas variáveis.

    \vspace{1em}
    O otimizador a ser utilizado durante os treinos das arquiteturas de \textit{CNN} é o \textit{Adam} (\textit{Adaptive Moment Estimation}), proposto por \textcite{kingma2017adammethodstochasticoptimization}.
    \vspace{1em}
    
    Demonstra ser mais eficiente devido ao pouco uso de memória durante o seu uso.
\end{frame}

\begin{frame}[fragile]{Otimizador}
    \begin{center}
        \resizebox{10.5cm}{!}{
        \centering
        \begin{minipage}{\linewidth}
            \begin{algorithm}[H]
                \caption{Algoritmo de otimização \textit{Adam} descrito por \textcite[tradução pelo autor]{kingma2017adammethodstochasticoptimization}}
                    \begin{algorithmic}
                        \Require $\alpha$\Comment{Taxa de aprendizado (\textit{learning rate})}
                        \Require $\beta_{1},\;\beta_{2}\in[0,1)$\Comment{Taxa de decaimento exponencial para a estimativa dos momentos}
                        \Require $f(\theta)$\Comment{Função objetiva estocástico com parâmetros $\theta$}
                        \Require $\theta_{0}$\Comment{Vetor de parâmetro inicial}
                            \While {$\theta_{0}$ não convergiu}
                            \State $t\leftarrow t+1$\Comment{Atualiza a iteração}
                            \State $g_t\leftarrow\nabla_{\theta}f_{t}(\theta_{t-1})$\Comment{Obtenção dos gradientes}
                            \State $m_t\leftarrow\beta_1\cdot m_{t-1}+(1-\beta_{1})\cdot g_{t}^{1}$\Comment{Atualiza o primeiro momento estimado}
                            \State $v_t\leftarrow\beta_{2}\cdot v_{t-1}+(1-\beta_2)\cdot g_{t}^{2}$\Comment{Atualiza o segundo momento estimado}
                            \State $\hat{m}_t\leftarrow\dfrac{m_{t}}{1-\beta^{t}_{1}}$\Comment{Computa a correção da estimativa do primeiro momento}
                            \State $\hat{v}_t\leftarrow\dfrac{v_{t}}{1-\beta^{t}_{2}}$\Comment{Computa a correção da estimativa do segundo momento}
                            \State $\theta_{t}\leftarrow\theta_{t-1}-\alpha\cdot\dfrac{\hat{m}_{t}}{\sqrt{\hat{v}_{t}}+\epsilon}$\Comment{Atualiza os parâmetros da função estocástica}
                            \EndWhile
                            \State \Return $\theta_{t}$ \Comment{Retorna o resultado dos parâmetros}
                    \end{algorithmic}
            \end{algorithm}
        \end{minipage}
        }
    \end{center}
    
\end{frame}